% Generated draft: requires mathematical and source review.
\documentclass[11pt]{article}
\usepackage[margin=25mm]{geometry}
\usepackage{fontspec}
\setmainfont{Latin Modern Roman}
\usepackage{amsmath,amssymb,mathtools,unicode-math}
\setmathfont{Latin Modern Math}
\usepackage{enumitem,needspace,xurl,hyperref}
\setlength{\parindent}{0pt}
\setlength{\parskip}{5pt}
\setlength{\emergencystretch}{3em}
\newcommand{\N}{\mathbb{N}}
\newcommand{\Z}{\mathbb{Z}}
\newcommand{\Q}{\mathbb{Q}}
\newcommand{\R}{\mathbb{R}}
\newcommand{\C}{\mathbb{C}}
\newcommand{\F}{\mathbb{F}}
\newcommand{\A}{\mathbb{A}}
\newcommand{\PP}{\mathbb{P}}
\newcommand{\E}{\mathbb{E}}
\newcommand{\eps}{\varepsilon}
\newcommand{\dd}{\,\mathrm{d}}
\begin{document}
\section*{UnsolvedMath ID: 30005797}
\textbf{Source problem code:} catalogue entry
\subsection{Definitions and mathematical statement}
The supplied passage concerns a singularity \(X\), a polytope denoted by \(P\), and a crepant resolution \(\widetilde X\) when such a resolution exists. It reports that the volume of \(P\) coincides with the Euler characteristic \(\chi(\widetilde X)\). The construction of \(P\), its volume normalization, the convention for Euler characteristic, and the hypotheses for this identification are not supplied.

A noncommutative crepant resolution is understood here in the sense of reference [4] in the report passage; that reference and its definition are not identified in the supplied record. The phrase ``noncommutative volume term'' denotes the replacement quantity associated with such resolutions mentioned in the record, without specifying a formula or asserting independence from the chosen resolution.

\textbf{Problem (interpretive formulation).} For singularities admitting noncommutative crepant resolutions within the class intended by the report, find an algebro-geometric construction or interpretation of the replacement quantity proposed by Gulotta in the report's reference [1]. This quantity is intended to replace the Euler characteristic of a crepant resolution in an analogue of the non-symmetric Mahler inequality. Explain its geometric meaning and its relationship to the volume term in that analogue.

The supplied excerpt does not specify the class of singularities, the replacement quantity, or a legible version of inequality (1). Consequently, it does not determine a precise numerical inequality or a fully quantified conjecture. The task above retains the stated interpretive target subject to these unresolved qualifications.

\textbf{Supplied contextual claims.} The passage says that the Euler-characteristic version of the inequality fails even for non-abelian quotient singularities. It attributes the proposed replacement to Gulotta and reports that the adjusted inequality holds with equality for finite quotient singularities, with further verification for a range of examples including anticanonical cones over del Pezzo surfaces and singularities labelled \texttt{cAn}. These are supplied claims whose hypotheses and proofs are absent; no assertion covering every example in the latter classes is inferred.
\subsection{Short English statement}
Find an algebro-geometric meaning for the quantity associated with noncommutative crepant resolutions that replaces the Euler characteristic in the proposed singularity analogue of the non-symmetric Mahler inequality. The supplied record leaves the quantity, exact inequality, and admissible singularities unspecified, so this remains a qualified interpretive problem requiring mathematical review.
\subsection{Sources}
\begin{itemize}
\item Local UnsolvedMath registry, permanent ID \(30005797\), exact source problem code \texttt{catalogue entry}, title ``Algebro-Geometric Meaning of Noncommutative Volume Bounds.'' The supplied record labels the problem \texttt{partially\_solved} and the statement \texttt{corrected\_verified}; these are registry labels, not independent verification here.
\item \emph{K-Stability, Birational Geometry and Mirror Symmetry} (2024), Oberwolfach Reports, pp.\ 827--870, DOI \texttt{10.4171/owr/2024/14}. Supplied source and literature link: \url{https://doi.org/10.4171/owr/2024/14}. No page within this range identifying the passage is supplied.
\item The excerpt attributes the replacement proposal to Gulotta in its reference [1] and uses noncommutative crepant resolutions in the sense of its reference [4]. Bibliographic details for [1] and [4] are not supplied. No proposer is identified in the registry.
\item Literature assessment dated \texttt{2026-08-22}: the supplied curation describes the Mahler-type inequality as conjectural, reports the proposed noncommutative-crepant-resolution interpretation as verified for finite quotient singularities and several further examples, and describes the general interpretation as unresolved. This is supplied dated curation derived from the OWR record, not independently verified. The record explicitly calls for independent verification of time-sensitive status and cited claims.
\end{itemize}
\end{document}
